\section*{अध्याय \ref{ch:b3:exploration-seismology} --- भूकंपीय अन्वेषण}

\begin{solution}{exo:b3:exploration-seismology:1}
$Z = \rho v$: चॉक $2000\times2500 = \qty{5.00e6}{kg.m^{-2}.s^{-1}}$, शेल
$2300\times2300 = \qty{5.29e6}{kg.m^{-2}.s^{-1}}$, चूना-पत्थर $2600\times4500 =
\qty{11.70e6}{kg.m^{-2}.s^{-1}}$। चॉक--शेल: $R = (5.29 - 5.00)/(5.29 + 5.00)
= 0.29/10.29 = 0.0282$; शेल--चूना-पत्थर: $R = (11.70 - 5.29)/(11.70 + 5.29) =
6.41/16.99 = 0.377$। परावर्तित ऊर्जा $R^2 = 0.142$ है: चूना-पत्थर की ऊपरी सतह
आपतित ऊर्जा का \qty{14.2}{\percent} वापस भेजती है, जबकि चॉक--शेल
अंतरापृष्ठ \qty{0.1}{\percent} से कम वापस भेजता है।
\end{solution}

\begin{solution}{exo:b3:exploration-seismology:2}
$t_0 = 2\times800/2000 = \qty{0.800}{s}$; $t = \sqrt{0.800^2 + (1000/2000)^2} =
\sqrt{0.89} = \qty{0.9434}{s}$; मूवआउट \qty{0.1434}{s} है। सन्निकटन
$1000^2/(2\times2000^2\times0.800) = \qty{0.1563}{s}$ देता है, जो
\qty{9}{\percent} अधिक है: जब ऑफ़सेट गहराई से बड़ा हो, तब $x \ll vt_0$
अब लागू नहीं होता।
\end{solution}

\begin{solution}{exo:b3:exploration-seismology:3}
मध्यबिंदु ग्राही-अंतराल के आधे के बराबर खिसकते हैं: CMP अंतराल
$12.5/2 = \qty{6.25}{m}$ है। आच्छादन गुणक $240\times12.5/(2\times25) = 60$ है।
120 के आच्छादन गुणक के लिए हर $240\times12.5/(2\times120) = \qty{12.5}{m}$ पर एक विस्फोट चाहिए।
\end{solution}

\begin{solution}{exo:b3:exploration-seismology:4}
$\lambda = v/f$: बलुआ पत्थर में \qty{100}{m}, $\lambda/4 = \qty{25}{m}$;
चूना-पत्थर में \qty{150}{m}, $\lambda/4 = \qty{37.5}{m}$। $\lambda = 2800/30 = \qty{93.3}{m}$ और
$z = \qty{1000}{m}$ के साथ,
$r_F \approx \sqrt{\lambda z/2} = \sqrt{46\,667} = \qty{216}{m}$ (\qty{217}{m}
पद $\lambda^2/16$ सहित)।
\end{solution}

\begin{solution}{exo:b3:exploration-seismology:5}
डिक्स: $v^2 = (2800^2\times0.700 - 2400^2\times0.500)/0.200 = (5\,488\,000 -
2\,880\,000)/0.200 = \qty{1.304e7}{m^2.s^{-2}}$, अतः $v = \qty{3611}{m.s^{-1}}$।
परत को एक तरफ़ा \qty{0.100}{s} में पार किया जाता है: $h = 3611\times0.100 =
\qty{361}{m}$।
\end{solution}

\begin{solution}{exo:b3:exploration-seismology:6}
(क) मान लीजिए $d$ मध्यबिंदु M से परावर्तक की दूरी है, ताकि
$t_0 = 2d/v$। स्रोत S से ग्राही G तक परावर्तित किरण उतनी ही लंबी है जितनी
परावर्तक में S के प्रतिबिंब S$'$ से G तक की सीधी रेखा। M को मूलबिंदु पर रखा जाता है और रेखा नति की दिशा में ली जाती है;
S$'$ अभिलंब के अनुदिश S से $2d_S$ दूरी पर है, जहाँ
$d_S = d - (x/2)\sin\varphi$ है। S$'$G का वर्ग है
$x^2 + 4d_S^2 + 4d_Sx\sin\varphi = x^2 + 4d^2 - x^2\sin^2\varphi$, सभी पद
$dx$ वाले निरस्त हो जाते हैं, अतः $v^2t^2 = 4d^2 + x^2\cos^2\varphi$ और
$t^2 = t_0^2 + x^2\cos^2\varphi/v^2$। (ख) गैदर एक ऐसे अतिपरवलय का अनुसरण करता है जिसका वेग
$V = v/\cos\varphi = 2000/\cos\ang{12} = \qty{2045}{m.s^{-1}}$ है: नति के कारण
\omterm{def:b3:exploration-seismology:velocities}{स्टैकिंग वेग} अत्यधिक ऊँचा हो जाता है, यहाँ \qty{2.2}{\percent} तक।
\end{solution}

\begin{solution}{exo:b3:exploration-seismology:7}
$\sqrt{60} = 7.75$ से। सिग्नल-से-शोर अनुपात दोगुना करने के लिए चार गुना
ट्रेस चाहिए: आच्छादन गुणक 240।
\end{solution}

\begin{solution}{exo:b3:exploration-seismology:8}
आलेखित गहराई $vt/2$ $\qty{1.5}{km.s^{-1}}\times\qty{0.40}{s}
= \qty{0.60}{km}$ प्रति किलोमीटर बदलती है: $\tan\theta_a = 0.60$, $\theta_a = \ang{31.0}$।
शून्य-ऑफ़सेट प्रमेय से $\sin\theta = 0.60$: $\theta = \ang{36.9}$।
\end{solution}

\begin{solution}{exo:b3:exploration-seismology:9}
हर स्पाइक उस पर केंद्रित $w$ की एक प्रति रखता है: नमूने 0--2 पर $0.3\,w$ से
$(-0.15, 0.3, -0.15)$ मिलता है, और नमूने 2--4 पर $-0.2\,w$ से $(0.1, -0.2, 0.1)$।
इनका योग $(0, -0.15, 0.3, -0.05, -0.2, 0.1, 0)$ है, जो श्रेणी से एक नमूना पहले से
एक नमूना बाद तक लिखा गया है। परावर्तन नमूने 1 और 3 पर शिखर बनाते हैं, अपने वास्तविक चिह्नों के साथ;
वे नमूने 2 पर अतिव्यापित होते हैं, जहाँ $-0.15 + 0.1 = -0.05$ है:
थोड़ा व्यतिकरण, पर दो अलग शिखर।
\end{solution}

\begin{solution}{exo:b3:exploration-seismology:10}
$\arcsin(1800/2900) = \ang{38.37}$ और $\arcsin(1800/4400) = \ang{24.15}$, अतः
$i_c = (38.37 + 24.15)/2 = \ang{31.26}$ और $\theta = (38.37 - 24.15)/2 =
\ang{7.11}$; $v_2 = 1800/\sin\ang{31.26} = \qty{3469}{m.s^{-1}}$। हरात्मक माध्य
$2\times2900\times4400/(2900 + 4400) = \qty{3496}{m.s^{-1}}$ इससे
\qty{0.8}{\percent} अधिक है; इसे $\cos\theta = 0.9923$ से गुणा करने पर
\qty{3469}{m.s^{-1}} मिलता है, जो सामान्य लघु-नति सन्निकटन
$v_2 \approx \cos\theta\,2v_dv_u/(v_d + v_u)$ है।
\end{solution}

\begin{solution}{exo:b3:exploration-seismology:11}
(क) संयुक्त संकेत $R[w(t) - w(t - \Delta t)] \approx R\,\Delta t\,w'(t)$ है,
जो प्रथम-कोटि टेलर विस्तार से प्राप्त होता है। (ख) इसका आयाम
$\Delta t$ के समानुपाती है, इसलिए \omterm{def:b3:exploration-seismology:resolution}{समस्वरण मोटाई} से पतला संस्तर भी एक संकेत लौटाता है,
जिसका आकार $w'$ का होता है, पर उसके शीर्ष और आधार को दो घटनाओं के रूप में नहीं चुना जा सकता:
वह पहचाना जाता है, पर विभेदित नहीं होता। यदि $R$ और तरंगिका ज्ञात हों, जैसे किसी कुएँ से,
तो आयाम से $\Delta t$ और इसलिए मोटाई प्राप्त होती है। (ग) आयाम मोटाई के
समानुपाती है: अनुपात $5/10 = 1/2$ है।
\end{solution}

\begin{solution}{exo:b3:exploration-seismology:12}
(क) $t(x) = (2/v)\sqrt{z^2 + x^2}$: $t(0) = 2\times600/3000 = \qty{0.400}{s}$ और
$t(300) = (2/3000)\sqrt{600^2 + 300^2} = \qty{0.4472}{s}$। (ख) $t^2 = t_0^2 +
x^2/(v/2)^2$, गति $v/2$ वाला NMO रूप है। उसी शीर्ष-काल पर परावर्तन,
शून्य-ऑफ़सेट खंड पर कोई मूवआउट नहीं दिखाता; गैदर में
उसी $t_0$ के NMO अतिपरवलय की तुलना में विवर्तन का मूवआउट
$x^2/(2(v/2)^2t_0)$ चार गुना बड़ा है, क्योंकि समय
प्रकीर्णक से दूरी पर दो बार निर्भर करता है, नीचे जाने और वापस आने में। (ग) माइग्रेशन
इस अतिपरवलय के अनुदिश खंड का योग करता है और ऊर्जा को उसके शीर्ष पर रखता है: विवर्तन
$x = 0$, $z = \qty{600}{m}$ पर एक बिंदु पर संकुचित हो जाता है।
\end{solution}

\begin{solution}{pb:b3:exploration-seismology:1}
\textbf{1.} W से, $400/(486.1 - 425.2) = 400/0.0609 = \qty{6568}{m.s^{-1}}$;
E से, $400/(486.1 - 372.4) = 400/0.1137 = \qty{3518}{m.s^{-1}}$। हेड
तरंग W से आने पर तेज़ लगती है: W नति का निचला सिरा है, और उसके विस्फोट ऊपर-नति की ओर यात्रा करते हैं; अतः
अपवर्तक W से E की ओर उठता है।

\textbf{2.} $\arcsin(2400/3518) = \ang{43.02}$, $\arcsin(2400/6568) =
\ang{21.43}$: $i_c = \ang{32.22}$, $\theta = \ang{10.79}$।

\textbf{3.} $v_2 = 2400/\sin\ang{32.22} = \qty{4501}{m.s^{-1}}$: यह चूना-पत्थर की परास में है,
जो जुरासिक चूना-पत्थर का शीर्ष है।

\textbf{4.} \qty{1200}{m} पर अंतःखंड $t - x/v_{\text{आभासी}}$ से प्राप्त होते हैं:
$t_{i,W} = 0.4861 - 1200/6568 = \qty{0.3034}{s}$ और $t_{i,E} = 0.4861 -
1200/3518 = \qty{0.1450}{s}$। $\cos i_c = 0.8460$ मानकर, $h = v_1t_i/(2\cos
i_c)$: $h_W = 2400\times0.3034/1.692 = \qty{430.4}{m}$ और $h_E =
\qty{205.7}{m}$।

\textbf{5.} $\cos\theta = 0.9823$ से भाग देने पर: W के नीचे \qty{438.1}{m} और
\qty{209.4}{m}, E के नीचे। फिर $(438.1 - 209.4)/1200 = 0.1906 =
\tan\ang{10.79}$, जो प्रश्न 2 की नति है।

\textbf{6.} शेल कोई हेड तरंग नहीं लौटाता: वह अदृश्य है, और सर्वेक्षण
उसे चॉक के साथ गिनता है। यहाँ $v_1 = \qty{2.40}{km.s^{-1}}$ दोनों का
माध्य मानकर दिया गया है; यदि चॉक की प्रत्यक्ष-तरंग चाल, \qty{2.50}{km.s^{-1}},
पूरे अधिभार के लिए प्रयुक्त की गई होती, तो हर गहराई बहुत अधिक निकलती।

\textbf{7.} $V = 600/\sqrt{0.5437^2 - 0.4878^2} = 600/0.2401 =
\qty{2499}{m.s^{-1}}$ और $V = 600/\sqrt{0.6240^2 - 0.5748^2} =
\qty{2471}{m.s^{-1}}$। दूसरा परिकलन चॉक को उसके नीचे के धीमे शेल के साथ
औसत करता है: माध्य घट जाता है।

\textbf{8.} \qty{2811}{m.s^{-1}} और \qty{2829}{m.s^{-1}}।

\textbf{9.} परत 1: \qty{2499}{m.s^{-1}}। परत 2: $v^2 = (2471^2\times0.5748
- 2499^2\times0.4878)/0.0870$, $v = \qty{2308}{m.s^{-1}}$। परत 3:
$v^2 = (2811^2\times0.6548 - 2471^2\times0.5748)/0.0800$,
$v = \qty{4561}{m.s^{-1}}$। परत 4: $v^2 = (2829^2\times0.7348 -
2811^2\times0.6548)/0.0800$, $v = \qty{2972}{m.s^{-1}}$।

\textbf{10.} चॉक (\qty{2499}{m.s^{-1}}), शेल (\qty{2308}{m.s^{-1}}),
चूना-पत्थर (\qty{4561}{m.s^{-1}}), बलुआ पत्थर (\qty{2972}{m.s^{-1}}), प्रत्येक अपनी
प्रकाशित परास के भीतर है। परावर्तन ये हैं: चॉक--शेल सीमा-पृष्ठ, चूना-पत्थर का शीर्ष,
चूना-पत्थर--बलुआ पत्थर सीमा-पृष्ठ और
शिस्ट आधार-शैल के शीर्ष पर की असंगति।

\textbf{11.} चॉक और शेल की ध्वनिक प्रतिबाधा लगभग एक-सी है: $R = 0.028$
(अभ्यास 1), जबकि चूना-पत्थर के शीर्ष पर यह 0.38 है।

\textbf{12.} वेग $V_{\text{वर्ग-माध्य-मूल}}$ वाला अतिपरवलय केवल ऑफ़सेट में द्वितीय कोटि तक
लागू होता है; \qty{2}{km} के ऑफ़सेट पर, जो सबसे गहरे परावर्तक की गहराई का
दोगुना है, समय उससे विचलित हो जाते हैं, और पूरे प्रसार पर फिट किया गया
अतिपरवलय गलत वेग देगा, और डिक्स अंतराल तो और भी बुरे होंगे। दूर के ऑफ़सेट
सुधार के समय तरंगिकाओं को भी खींच देते हैं।

\textbf{13.} परतों में द्वि-मार्गी समय: \qty{0.4878}{s}, \qty{0.0870}{s},
\qty{0.0800}{s}, \qty{0.0800}{s}। मोटाइयाँ $v\,\Delta t/2$: \qty{609.5}{m},
\qty{100.4}{m}, \qty{182.4}{m}, \qty{118.9}{m}।

\textbf{14.} ज़मीन के नीचे गहराइयाँ: \qty{609.5}{m}, \qty{709.9}{m},
\qty{892.3}{m}, \qty{1011.2}{m}। शेल, चूना-पत्थर और बलुआ पत्थर
अपनी मोटाइयों से \qty{0.4}{m}, \qty{2.4}{m} और \qty{1.1}{m} के भीतर आते हैं, त्रुटि
तेज़ और पतले चूना-पत्थर में बढ़ती है, जहाँ समय की छोटी-सी त्रुटि
सबसे अधिक महँगी पड़ती है।

\textbf{15.} $\bar v = 2\times1011.2/0.7348 = \qty{2752}{m.s^{-1}}$। RMS
वेग प्रत्येक परत को $v^2$ से भारित करता है और कभी भी समय-भारित
माध्य से कम नहीं होता: $2829\times0.7348/2 = \qty{1039}{m}$, जो \qty{28}{m} अधिक गहरा है।

\textbf{16.} $\lambda = 2308/30 = \qty{76.9}{m}$, $\lambda/4 = \qty{19.2}{m}$;
\omterm{def:b3:exploration-seismology:wavelet}{रिकर तरंगिका} का अनुमान $2308/(5.2\times30) = \qty{14.8}{m}$ है। शेल
\qty{100}{m} मोटा है, अपनी \omterm{def:b3:exploration-seismology:resolution}{समस्वरण मोटाई} का पाँच गुना: उसका शीर्ष और तल
विभेदित हैं। चतुर्थांश-तरंगदैर्घ्य नियम से यह समस्वरित होता, $f = v/(4b) =
2308/400 = \qty{5.77}{Hz}$ आवृत्ति पर।

\textbf{17.} $\lambda = 2829/30 = \qty{94.3}{m}$, $z = \qty{1011}{m}$:
$r_F \approx \sqrt{94.3\times1011/2} = \qty{218}{m}$ (\qty{220}{m}, जब
पद $\lambda^2/16$ जोड़ा जाए)।

\textbf{18.} $t(x)^2 = t_0^2 + (2x/v)^2$, अतः $v = 2x/\sqrt{t^2 - t_0^2} =
400/\sqrt{0.5732^2 - 0.550^2} = 400/0.1614 = \qty{2478}{m.s^{-1}}$।

\textbf{19.} $z = vt_0/2 = 2478\times0.550/2 = \qty{681}{m}$।

\textbf{20.} आलेखित गहराई $vt/2$ प्रति इकाई लंबाई पर $1.24\times0.167 = 0.207$ से
बदलती है: $\theta_a = \arctan0.207 = \ang{11.70}$। वास्तविक नति के लिए
$\sin\theta = 0.207$ है: $\theta = \ang{11.95}$, वलन-भुजा की नति।

\textbf{21.} $d = 2480\times0.655/2 = \qty{812}{m}$; माइग्रेशन परावर्तन को
$d\sin\theta = 812\times0.207 = \qty{168}{m}$ ऊपर-नति की ओर, अपनति के
शिखर की ओर, और \qty{812}{m} से ऊपर उठाकर $d\cos\theta =
\qty{795}{m}$ तक ले जाता है।

\textbf{22.} $2500\times0.160/2 = \qty{200}{m}$।

\textbf{23.} प्रश्न 13 और 14 से पता चलता है कि आधार-शैल की असंगति
\textbf{\qty[text-series-to-math=true]{1011}{m}} ज़मीन के नीचे CMP S पर स्थित है।
\end{solution}
